\documentclass[12pt,a4paper]{article}

\usepackage[utf8]{inputenc}
\usepackage[T1]{fontenc}
\usepackage[spanish,es-nodecimaldot]{babel}
\usepackage{geometry}
\usepackage{graphicx}
\usepackage{float}
\usepackage{xcolor}
\usepackage{amsmath}
\usepackage{listings}
\usepackage{fancyhdr}

\geometry{margin=2.5cm}
\linespread{1.15}
\setlength{\headheight}{14.5pt}
\setlength{\emergencystretch}{2em}

\definecolor{codebg}{RGB}{248,248,248}
\definecolor{codeframe}{RGB}{205,210,220}
\definecolor{codekey}{RGB}{28,80,170}
\definecolor{codecomment}{RGB}{34,125,72}
\definecolor{codestring}{RGB}{170,55,45}
\definecolor{codenum}{RGB}{120,120,120}
\definecolor{unsaacblue}{RGB}{31,82,150}

\lstdefinestyle{cxxcolor}{
  language=C++,
  basicstyle=\ttfamily\small,
  backgroundcolor=\color{codebg},
  frame=single,
  rulecolor=\color{codeframe},
  keywordstyle=\color{codekey}\bfseries,
  commentstyle=\color{codecomment}\itshape,
  stringstyle=\color{codestring},
  numberstyle=\tiny\color{codenum},
  numbers=left,
  stepnumber=1,
  numbersep=8pt,
  breaklines=true,
  showstringspaces=false,
  columns=fullflexible,
  keepspaces=true,
  tabsize=2
}

\pagestyle{fancy}
\fancyhf{}
\rhead{Pr\'actica 11}
\lhead{Transformaciones Avanzadas 3D}
\cfoot{\thepage}

\begin{document}

\begin{titlepage}
	\centering
	
\includegraphics[width=3.2cm]{\detokenize{../../logo_unsaac.jpeg}}\par
	\vspace{0.8cm}
	{\Large \textbf{UNIVERSIDAD NACIONAL DE SAN ANTONIO ABAD DEL CUSCO}}\\[0.2cm]
	{\large Departamento Acad\'emico de Inform\'atica}\\[0.1cm]
	{\large Escuela Profesional: Ing. Inform\'atica y de Sistemas}\\[1.3cm]

	{\Large \textbf{COMPUTACI\'ON GR\'AFICA I}}\\[0.4cm]
	{\Large \textbf{Pr\'actica N$^{\circ}$11}}\\[0.25cm]
	{\LARGE \textbf{TRANSFORMACIONES AVANZADAS EN 3D}}\\[0.25cm]
	{\Large \textbf{Reflexi\'on general y distorsi\'on din\'amica}}\\[1.2cm]

	\begin{flushleft}
	\textbf{Estudiante:} Efrain Vitorino Marin\\
	\textbf{C\'odigo:} 160337\\
	\textbf{Profesor:} Dr. Hans Harley Ccacyahuillca Bejar\\
	\textbf{Repositorio:} \texttt{https://github.com/lovito99/laboGrafica/tree/main/labo11}\\
	\end{flushleft}

	\vfill
	{\large Cusco -- Per\'u\\2026}
\end{titlepage}

\section*{Resumen}
En esta pr\'actica se desarrollaron tres transformaciones avanzadas en el espacio tridimensional usando OpenGL moderno. Se implement\'o la reflexi\'on de una pir\'amide sobre el plano arbitrario $3x+5y+6z=0$, la reflexi\'on estricta sobre el plano infinito $Z=5$ y una distorsi\'on din\'amica cruzada aplicada sobre la pir\'amide reflejada.

\section*{Objetivos}
\begin{itemize}
	\item Implementar reflexiones tridimensionales mediante matrices homog\'eneas de $4\times4$.
	\item Aplicar una reflexi\'on general respecto a un plano arbitrario que pasa por el origen.
	\item Deducir el t\'ermino de traslaci\'on para reflejar respecto al plano $Z=5$.
	\item Construir una matriz de cizallamiento 3D con factores dependientes del tiempo.
	\item Visualizar los resultados con OpenGL moderno, GLFW, GLEW y GLM.
\end{itemize}

\section*{Librer\'ias usadas}
La implementaci\'on mantiene el estilo de la plantilla del PDF, usando OpenGL moderno, shaders y matrices enviadas como variables uniformes.

\begin{lstlisting}[style=cxxcolor,caption={Librerias principales usadas.}]
#include <GL/glew.h>
#include <GLFW/glfw3.h>
#include <glm/glm.hpp>
#include <glm/gtc/matrix_transform.hpp>
#include <glm/gtc/type_ptr.hpp>
\end{lstlisting}

\section*{Shaders}
El \textit{vertex shader} recibe la posici\'on de cada v\'ertice y aplica las matrices \texttt{projection}, \texttt{view} y \texttt{transformMatrix}. El \textit{fragment shader} colorea cada objeto mediante una variable uniforme.

\begin{lstlisting}[style=cxxcolor,caption={Vertex shader usado en la practica.}]
const char* const shaderVert = R"glsl(
#version 330 core
layout (location = 0) in vec3 aPos;
uniform mat4 transformMatrix;
uniform mat4 view;
uniform mat4 projection;

void main()
{
    gl_Position = projection * view * transformMatrix * vec4(aPos, 1.0);
}
)glsl";
\end{lstlisting}

\begin{lstlisting}[style=cxxcolor,caption={Fragment shader usado para colorear la piramide y los ejes.}]
const char* const shaderFrag = R"glsl(
#version 330 core
out vec4 FragColor;
uniform vec3 objectColor;

void main()
{
    FragColor = vec4(objectColor, 1.0);
}
)glsl";
\end{lstlisting}

\section*{Geometr\'ia de la escena}
La escena contiene ejes coordenados y una pir\'amide definida con cuatro v\'ertices. La pir\'amide se dibuja mediante aristas para poder comparar con claridad la figura original y la figura transformada.

\begin{lstlisting}[style=cxxcolor,caption={Vertices de la piramide usada en las tres actividades.}]
const float pyramidVertices[] = {
    40.0f, 20.0f, 40.0f,
    60.0f, 30.0f, 20.0f,
    40.0f, 50.0f, 30.0f,
    20.0f, 30.0f, 20.0f
};

const unsigned int pyramidEdges[] = {
    0, 1, 1, 2, 2, 0,
    0, 3, 1, 3, 2, 3
};
\end{lstlisting}

\section*{C\'amara y proyecci\'on}
Para observar las transformaciones en 3D se usa una proyecci\'on ortogr\'afica y una c\'amara construida con \texttt{glm::lookAt}.

\begin{lstlisting}[style=cxxcolor,caption={Envio de matrices de camara al shader.}]
inline void sendCamera(GLuint program)
{
    const glm::mat4 projection = glm::ortho(
        -120.0f, 120.0f,
        -100.0f, 130.0f,
        -250.0f, 250.0f);

    const glm::mat4 view = glm::lookAt(
        glm::vec3(95.0f, 80.0f, 145.0f),
        glm::vec3(25.0f, 25.0f, 25.0f),
        glm::vec3(0.0f, 1.0f, 0.0f));

    glUniformMatrix4fv(glGetUniformLocation(program, "projection"),
                       1, GL_FALSE, glm::value_ptr(projection));
    glUniformMatrix4fv(glGetUniformLocation(program, "view"),
                       1, GL_FALSE, glm::value_ptr(view));
}
\end{lstlisting}

\section*{Actividad A: Reflexi\'on sobre $3x+5y+6z=0$}
El plano pasa por el origen y su vector normal es:

\[
\vec{n}=(3,5,6)
\]

Para reflejar un punto $p$ respecto a un plano que pasa por el origen se usa:

\[
p' = p - 2\vec{n}\frac{\vec{n}\cdot p}{\vec{n}\cdot \vec{n}}
\]

En forma matricial:

\[
R = I - 2\frac{\vec{n}\vec{n}^{T}}{\vec{n}\cdot\vec{n}}
\]

Como:

\[
\vec{n}\cdot\vec{n}=3^2+5^2+6^2=70
\]

la matriz queda:

\[
R =
\begin{pmatrix}
1-\frac{18}{70} & -\frac{30}{70} & -\frac{36}{70} & 0\\
-\frac{30}{70} & 1-\frac{50}{70} & -\frac{60}{70} & 0\\
-\frac{36}{70} & -\frac{60}{70} & 1-\frac{72}{70} & 0\\
0 & 0 & 0 & 1
\end{pmatrix}
\]

\begin{lstlisting}[style=cxxcolor,caption={Construccion de la reflexion sobre un plano que pasa por el origen.}]
inline void createReflectionPlaneThroughOrigin(float* mat, float a, float b, float c)
{
    createIdentityMatrix(mat);

    const float denominator = a * a + b * b + c * c;
    if (denominator == 0.0f) {
        return;
    }

    const float n[3] = {a, b, c};
    for (int row = 0; row < 3; ++row) {
        for (int col = 0; col < 3; ++col) {
            const float identity = row == col ? 1.0f : 0.0f;
            mat[col * 4 + row] =
                identity - (2.0f * n[row] * n[col]) / denominator;
        }
    }
}
\end{lstlisting}

\begin{lstlisting}[style=cxxcolor,caption={Uso de la matriz en la actividad A.}]
float identity[16];
float reflection[16];
createIdentityMatrix(identity);
createReflectionPlaneThroughOrigin(reflection, 3.0f, 5.0f, 6.0f);

drawPyramid(buffers, transformLoc, colorLoc, identity, 0.90f, 0.12f, 0.10f);
drawPyramid(buffers, transformLoc, colorLoc, reflection, 0.05f, 0.52f, 0.85f);
\end{lstlisting}

\begin{figure}[H]
	\centering
	\includegraphics[width=0.86\linewidth]{\detokenize{../refliexion plano.png}}
	\caption{Reflexi\'on de la pir\'amide sobre el plano $3x+5y+6z=0$.}
\end{figure}

\section*{Actividad B: Reflexi\'on estricta sobre $Z=5$}
Para reflejar respecto al plano infinito $Z=5$, se traslada el plano hacia el origen, se refleja sobre $Z=0$ y se retorna a la posici\'on original:

\[
M = T(0,0,5)\,R_{xy}\,T(0,0,-5)
\]

La transformaci\'on resultante es:

\[
x'=x,\qquad y'=y,\qquad z'=10-z
\]

La matriz homog\'enea queda:

\[
M =
\begin{pmatrix}
1 & 0 & 0 & 0\\
0 & 1 & 0 & 0\\
0 & 0 & -1 & 10\\
0 & 0 & 0 & 1
\end{pmatrix}
\]

El valor $10$ corresponde al t\'ermino de traslaci\'on acoplado $2\cdot5$.

\begin{lstlisting}[style=cxxcolor,caption={Reflexion sobre el plano Z igual a una constante.}]
inline void createReflectionZEquals(float* mat, float zPlane)
{
    createIdentityMatrix(mat);
    mat[10] = -1.0f;
    mat[14] = 2.0f * zPlane;
}
\end{lstlisting}

\begin{lstlisting}[style=cxxcolor,caption={Uso de la reflexion estricta sobre Z=5.}]
float identity[16];
float reflectionZ5[16];
createIdentityMatrix(identity);
createReflectionZEquals(reflectionZ5, 5.0f);

drawPyramid(buffers, transformLoc, colorLoc, identity, 0.90f, 0.12f, 0.10f);
drawPyramid(buffers, transformLoc, colorLoc, reflectionZ5, 0.06f, 0.58f, 0.20f);
\end{lstlisting}

\begin{figure}[H]
	\centering
	\includegraphics[width=0.86\linewidth]{\detokenize{../reflexion estricta.png}}
	\caption{Reflexi\'on geom\'etrica estricta de la pir\'amide sobre el plano $Z=5$.}
\end{figure}

\section*{Actividad C: Distorsi\'on din\'amica cruzada}
La actividad final aplica cizallamiento 3D sobre la pir\'amide reflejada. La matriz general usada es:

\[
Sh =
\begin{pmatrix}
1 & shxy & shxz & 0\\
shyx & 1 & shyz & 0\\
shzx & shzy & 1 & 0\\
0 & 0 & 0 & 1
\end{pmatrix}
\]

En esta pr\'actica se solicita:

\[
shxy=\sin(t), \qquad shyz=\cos(t)
\]

Por tanto:

\[
x'=x+\sin(t)y
\]

\[
y'=y+\cos(t)z
\]

\[
z'=z
\]

\begin{lstlisting}[style=cxxcolor,caption={Construccion de matriz de shearing 3D.}]
inline void createShearing3D(float* mat, float shxy, float shxz,
                             float shyx, float shyz,
                             float shzx, float shzy)
{
    createIdentityMatrix(mat);
    mat[4] = shxy;
    mat[8] = shxz;
    mat[1] = shyx;
    mat[9] = shyz;
    mat[2] = shzx;
    mat[6] = shzy;
}
\end{lstlisting}

\begin{lstlisting}[style=cxxcolor,caption={Distorsion dinamica aplicada sobre la piramide reflejada.}]
const float time = static_cast<float>(glfwGetTime());
float shear[16];
float dynamicTransform[16];

createShearing3D(
    shear,
    std::sin(time), 0.0f,
    0.0f, std::cos(time),
    0.0f, 0.0f);

multiplyMatrices(shear, reflectionZ5, dynamicTransform);

drawPyramid(buffers, transformLoc, colorLoc, identity, 0.72f, 0.72f, 0.72f);
drawPyramid(buffers, transformLoc, colorLoc, reflectionZ5, 0.06f, 0.58f, 0.20f);
drawPyramid(buffers, transformLoc, colorLoc, dynamicTransform, 0.88f, 0.28f, 0.06f);
\end{lstlisting}

\begin{figure}[H]
	\centering
	\includegraphics[width=0.86\linewidth]{\detokenize{../distorcion dinamica.png}}
	\caption{Distorsi\'on din\'amica cruzada aplicada sobre la pir\'amide reflejada.}
\end{figure}

\section*{Compilaci\'on y ejecuci\'on}
Los tres programas se compilan con CMake:

\begin{lstlisting}[style=cxxcolor,caption={Comandos de compilacion.}]
cd ~/Documentos/GitHub/laboGrafica/labo11
mkdir -p build
cd build
cmake .. -G "Unix Makefiles"
make
\end{lstlisting}

Luego se ejecutan:

\begin{lstlisting}[style=cxxcolor,caption={Comandos de ejecucion.}]
./actARefPlano
./actBRefZ5
./actCShearDyn
\end{lstlisting}

\section*{Conclusiones}
\begin{itemize}
	\item La reflexi\'on sobre el plano $3x+5y+6z=0$ se obtuvo con una matriz general basada en el vector normal del plano.
	\item La reflexi\'on sobre $Z=5$ requiere incorporar el t\'ermino de traslaci\'on $10$ dentro de la matriz homog\'enea.
	\item La distorsi\'on din\'amica cruzada se logr\'o modificando simult\'aneamente las componentes $X$ e $Y$ con factores dependientes del tiempo.
	\item El uso de shaders y matrices uniformes permite separar la geometr\'ia original de las transformaciones aplicadas.
\end{itemize}

\end{document}
